\section{Near-optimal factors, resilience and many prime Hamilton cycles}
\label{sec:regularfactors}
The weighted rounding theorem permits very nonuniform edge weights.
In the integer prime-sum application there is additional uniformity:
ordinary, unweighted degrees are asymptotically equal.
This allows an integral regular-factor argument, proved here, which
strengthens the arithmetic conclusion substantially.
It is an alternative application mechanism; adaptive rounding is
not claimed to be necessary for this particular uniform graph.

\begin{lemma}[Regular factors from near regularity and cuts]
\label{lem:regularfactor}
Let \(G\) be a simple balanced bipartite graph on \(N=2m\)
vertices. Suppose, for real \(D>0\), \(0<g\leq1\) and
\(0\leq\varepsilon<g/2\), that
\[
 (1-\varepsilon)D\leq\deg_G(v)\leq(1+\varepsilon)D
 \quad\hbox{for every }v,
\]
and
\[
 e_G(W,W^c)\geq gD\min\{|W|,N-|W|\}
 \quad\hbox{for every }W\subseteq V(G).
 \tag{21}
\]
Then \(G\) contains an \(r\)-regular spanning subgraph whenever
\(r\) is a nonnegative integer with
\[
 r\leq(1-2\varepsilon/g)D.                                  \tag{22}
\]
In particular one may take the floor of the right side.
Every such factor \(H\) satisfies the finite operator bound
\[
 \|A_G-A_H\|\leq(1+\varepsilon)D-r.                          \tag{23}
\]
\end{lemma}
\begin{proof}
Write the parts as \(A,B\).
The finite matching criterion for an \(r\)-factor is
\[
 e_G(S,B\setminus T)\geq r(|S|-|T|)
 \quad(S\subseteq A,\ T\subseteq B).                         \tag{24}
\]
Here is a capacity-one construction proving sufficiency, including
\(r=0\). Form a bipartite matching gadget with left vertices
\(A\times\{1,\ldots,r\}\) and one \(L_e\) for each edge, and right
vertices \(B\times\{1,\ldots,r\}\) and one \(R_e\) for each edge.
Join an \(A\)-slot to \(R_e\) when \(e\) starts at that vertex,
join \(L_e\) to every slot of its \(B\)-endpoint, and join
\(L_e\) to \(R_e\).

For an arbitrary left set, let \(K\) be its slots, \(F\) its edge
vertices, \(S\) the vertices touched by \(K\), and \(T\) the
\(B\)-endpoints of \(F\). Its neighborhood contains all \(r|T|\)
slots over \(T\), all \(R_e\) with \(e\in F\), and all \(R_e\)
with \(e\in E(S,B\setminus T)\). These last two edge sets are
disjoint. Thus (24) gives
\[
 |N(K\cup F)|\geq r|T|+|F|+e_G(S,B\setminus T)
               \geq r|S|+|F|\geq |K|+|F|.
\]
Finite Hall matching applies. Both gadget parts have size
\(rm+|E(G)|\), so the matching is perfect. An edge \(e\) is
either matched internally \(L_eR_e\), or its two vertices are
matched to an \(A\)-slot and a \(B\)-slot. Select the latter
edges. Each slot is matched once and each edge has only one
private pair, so the selected graph is simple and has degree
exactly \(r\) at every vertex.

Only \(x=|S|-|T|>0\) matters.
By exchanging the two sides and replacing \((S,T)\) with
\((B\setminus T,A\setminus S)\), if necessary, we may assume
\[
 u=|S|+|T|\leq m.
\]
Both \(x\) and the edge count in (24) are preserved by this
replacement. Put
\[
 a=e_G(S,B\setminus T),\qquad b=e_G(A\setminus S,T).
\]
The full boundary of \(S\cup T\) is \(a+b\), so (21) and
the degree bounds give
\[
 a+b\geq gDu,\qquad a-b\geq Dx-\varepsilon Du.
\]
Since \(b\geq0\), the two simultaneous lower bounds are
\[
 a\geq D(x-\varepsilon u),\qquad
 a\geq \frac D2\bigl(x+(g-\varepsilon)u\bigr).                \tag{25}
\]
If \(x\leq(g-\varepsilon)u\), the second gives \(a\geq Dx\).
Otherwise the first gives
\[
 a\geq Dx\left(1-\frac{\varepsilon}{g-\varepsilon}\right)
       \geq Dx(1-2\varepsilon/g).
\]
Thus (24) follows from (22).
Finally \(A_G-A_H\) is a symmetric nonnegative matrix with
maximum row sum at most the right side of (23).
Its operator norm obeys the same bound.
\end{proof}

The complete finite lemma is formalized, with the actual edge set
and adjacency matrix. \texttt{HamiltonianFiniteHall} provides
the finite matching proof without rebuilding a dependency cache.
\texttt{HamiltonianRegularFactorHall} retains the complete slot
and private-edge gadget; \texttt{HamiltonianRegularFactorMatching}
extracts the simple factor and proves both exact degree equalities.
\texttt{HamiltonianFactorIncidence} proves the signed crossing
balance and its complementary degree bound.
\texttt{HamiltonianRegularFactorCriterion} combines these with
the same full cut and proves (22).
\texttt{HamiltonianRegularFactorOperator} proves the exact
adjacency row counts and (23), and assembles both conclusions
for the same factor. All thirty-four declarations in these six
owners pass the focused foundational-axiom audit. The actual
near-optimal prime-factor consumer, including the single family
and its native adjacency eigenvalues, is formalized below.

Let \(B_N\) denote the ordinary adjacency matrix of \(G_N\)
and put \(D_N^*=N/\log N\).
This scalar is distinct from the degree matrix \(D_N\)
in Lemma~\ref{lem:primeweights}.

\begin{lemma}[Unweighted spectral transfer]
\label{lem:unweightedtransfer}
As even \(N\) tends to infinity,
\[
 \max_i\left|\frac{\deg_{G_N}(i)}{D_N^*}-1\right|
       =O\left(\frac{\log\log N}{\log N}\right),\qquad
 \|B_N/D_N^*-Q_N\|
       =O\left(\frac{\log\log N}{\log N}\right).              \tag{26}
\]
There is a stochastic kernel \(P_N\) on the same support with
\[
 \|B_N/D_N^*-P_N\|=o(1),\qquad
 \|P_N|_{\{\one,s\}^{\perp}}\|\leq1/2+o(1).
 \tag{27}
\]
Consequently, for all sufficiently large \(N\),
\[
 e_{G_N}(W,W^c)\geq\frac15D_N^*
                  \min\{|W|,N-|W|\}                        \tag{28}
\]
for every vertex set \(W\).
\end{lemma}
\begin{proof}
The difference in (26) has entries
\[
 \frac{\log N-\log(i+j)}N
                    \one_{\,i\ne j,\ i+j\ {\rm prime}}.
\]
When \(i+j\geq N/(\log N)^4\), the absolute numerator is at
most \(4\log\log N+\log2\), and each row has at most
\(\pi(2N)=O(N/\log N)\) nonzero entries.
The remaining outputs number at most \(N/(\log N)^4\)
per row, each with absolute numerator at most \(\log(2N)\).
Thus the maximum absolute row sum of the difference is
\(O(\log\log N/\log N)\).
Symmetry proves the operator bound, while the row-sum version
and (20) prove the degree assertion.
The exact scaling already used in the proof of
Theorem~\ref{thm:filz} gives (27).

For the boundary estimate put \(w=|W|\) and
\(x=\one_W-(w/N)\one\).
The identity for the combinatorial Laplacian gives
\[
 e_{G_N}(W,W^c)
 =x^\top(\diag(B_N\one)-B_N)x.
\]
On \(\one^\perp\), the algebraically largest eigenvalue of
\(P_N\) is at most \(1/2+o(1)\); its parity eigenvalue is \(-1\).
Equations (26)--(27), with the degree lower bound, therefore give
\[
 e_{G_N}(W,W^c)
 \geq(1/2-o(1))D_N^*\|x\|_2^2
 =(1/2-o(1))D_N^*\frac{w(N-w)}N.
\]
Since \(w(N-w)/N\geq\frac12\min(w,N-w)\), this implies (28).
\end{proof}

\paragraph{Formalized complete unweighted arithmetic transfer.}
The actual constant-edge matrix
\[
 U_n(i,j)=\frac{\log n}{n}
      \one_{i\ne j,\ i+j\ {\rm prime}}
\]
is now defined in \texttt{HamiltonianUnweightedPrimeTransfer}.
The original loop and prime masks are preserved exactly.
For every positive integer \(n\) with \(\log n\geq64\), its
comparison with the true logarithmic prime matrix satisfies
\[
 \max_i\sum_j|U_n(i,j)-Q_n(i,j)|\leq E_n,\qquad
 \max_j\sum_i|U_n(i,j)-Q_n(i,j)|\leq E_n,\qquad
 \|U_n-Q_n\|\leq E_n,
\]
where the explicit error is
\[
 E_n=\frac{1+4\log4\,(2\log\log n+\log2)}{\log n}.
\]
\texttt{HamiltonianLogWeightComparison} proves the two
pointwise regimes and the complete finite packing estimate.
\texttt{HamiltonianPrimeLogTransfer} uses
\(K=\lfloor n/(\log n)^2\rfloor\), charges every small
output to its integer count and every large output to its
actual prime logarithmic weight, and consumes
\(\theta(2n)\leq2n\log4\).
No prime-count asymptotic or independent prime-pair estimate
is assumed in this transfer.
The injective physical row intervals, all endpoints and both
matrix directions are included before the separate operator
error is bounded.

\texttt{HamiltonianPrimeLogTransferDecay} proves the uniform
decay on an entire terminal interval, including
\[
 E_n\leq\frac{1+16\log4+4\log4\log2}{\sqrt{\log n}}.
\]
It also proves, for every natural \(b\), one \(D_b>0\)
and one starting order such that every later row satisfies
\[
 \left|\sum_jU_n(i,j)-1\right|
                   \leq E_n+D_b/(\log n)^b.
\]
Thus the logarithmic rate in (26) follows from an explicit
uniform inequality. The actual constant-edge spectral consumer
is formalized without an auxiliary grid:
\[
 \forall\varepsilon>0\ \exists n_0 \forall n\geq n_0:\qquad
 \|U_n-P_{0,n}\|\leq\tfrac12+\varepsilon .
\]
All twenty-four new declarations in these four maintained owners
pass the focused foundational-axiom audit.
The balanced odd/even cross-block identification and the actual
support-preserving stochastic repair are formalized, as detailed
in Section~\ref{sec:relativeprimecircles}. The same witness gives
the exact entropy maximum and the actual unoriented prime-circle
count, including the factorial normalization and the quantitative
Stirling conversion. Its two relative errors are respectively
at most \(m^{-3/5}\) and \(2m^{-3/5}\), \(m=N/2\).
\paragraph{Formalized sharp repair and the full prime-edge cut bound.}
\texttt{HamiltonianPrimeSharpScaling} retains the same positive
counting witness while making its perturbation arbitrarily small.
\texttt{HamiltonianPrimePhysicalScaling} transfers it back through
the exact odd/even bijection. For every \(\eta>0\), for every
sufficiently large \(m\), the actual physical matrix \(P_{2m}\)
is symmetric, nonnegative and exactly stochastic, with
\[
 P_{2m}(i,j)=0\ \Longleftrightarrow\ U_{2m}(i,j)=0,\qquad
 \|P_{2m}-U_{2m}\|\leq\eta,\qquad
 \|P_{2m}-P_{0,2m}\|\leq\tfrac12+\eta .
\]
\texttt{HamiltonianPrimeOrthogonalGap} proves the exact
constant/parity multiplication formula and the operator bound
on the subspace with both corresponding sums zero.
This completes the literal sharp-gap consumer in (27).

For (28), \texttt{HamiltonianSpectralCuts} and
\texttt{HamiltonianCenteredCuts} prove the full signed
Dirichlet identity and exact centered indicator norm.
For every centered cut vector \(x\),
\texttt{HamiltonianPrimeCuts} retains
\[
 x^\top P_{0,N}x=-\frac1N
       \left(\sum_i(-1)^i x_i\right)^2\leq0.
\]
The degree lower bound \(19/20\) and operator error at most
\(11/20\) give the complete finite bound
\[
 e_{U_N}(W,W^c)\geq\frac25\,\frac{|W|(N-|W|)}N
                       \geq\frac15\min\{|W|,N-|W|\}.
\]
The same owner defines the actual finite set of distinct
prime-sum edges crossing the cut, proves its exact weighted
cardinality identity, and concludes
\[
 e_{G_N}(W,W^c)\geq\frac{N}{5\log N}
                           \min\{|W|,N-|W|\}
\]
on an entire terminal interval, uniformly over every \(W\).
No cut or overlap is removed. All thirty-nine new declarations
in these six owners pass focused foundational-axiom checks;
the two audit batches are disjoint. Equations (26)--(28)
now have complete maintained arithmetic consumers.

\begin{theorem}[Near-optimal regular prime-sum factors]
\label{thm:nearlyoptimalfactor}
For every sufficiently large even \(N\), \(G_N\) has a simple
\(r_N\)-regular spanning subgraph \(J_N\) such that
\[
 r_N=(1-o(1))\frac N{\log N},\qquad
 \|A_{J_N}/r_N-P_N\|=o(1).
 \tag{29}
\]
In particular every eigenvalue other than its exact
\(r_N,-r_N\) has modulus at most \((1/2+o(1))r_N\).
The factor degree is asymptotically best possible.
\end{theorem}
\begin{proof}
Use Lemma~\ref{lem:regularfactor} with \(g=1/5\),
\(D=D_N^*\), and the vanishing degree defect from (26).
It supplies
\[
 r_N=\left\lfloor(1-10\varepsilon_N)D_N^*\right\rfloor
                      =(1-o(1))D_N^* .
\]
The discarded adjacency has norm \(o(D_N^*)\) by (23).
Together with (27) and \(r_N/D_N^*\to1\), this proves (29).
The vectors \(\one,s\) remain exact, so the asserted
nontrivial spectral bound follows on their common complement.
Any regular factor has degree at most the minimum degree of
\(G_N\), which is \((1+o(1))D_N^*\).
\end{proof}

The finite prime-factor consumer is now formalized.
\texttt{HamiltonianFactorCutGeometry} identifies both crossing
counts with the full cut boundary and transports it losslessly
through the odd/even vertex bijection.
\texttt{HamiltonianPrimeFactorInputs} defines the actual prime
edge set, proves both ordinary degree estimates with
\(D=2m/\log(2m)\), and transfers the cut coefficient \(1/5\)
to that same edge set.
\texttt{HamiltonianPrimeRegularFactor} proves that, for every
fixed \(0<\varepsilon<1/10\), there is \(M_0\) such that
every \(m\geq M_0\) and every integer
\(0\leq r\leq(1-10\varepsilon)D\) have a simple factor \(H\)
on the original integer vertices with
\[
 \deg_H(v)=r,\qquad E(H)\subseteq E(G_{2m}),\qquad
 \|A_{G_{2m}}-A_H\|\leq(1+\varepsilon)D-r.
\]
Its declarations check the original distinct-integer prime
support and exact physical degree and norm transfers.
All twenty-three declarations in these three owners pass
focused foundational-axiom checks.

The single shrinking-defect factor family and its complete
operator normalization are now formalized as well.
The simultaneous degree estimate supplies one fixed \(B>0\).
Write \(L=\log(2m)\) and \(D=2m/L\). For every sufficiently
large \(m\), the chosen factor has the fixed degree
\[
 r_m=\left\lfloor
       \left(1-\frac{10B}{\sqrt L}\right)D
     \right\rfloor,\qquad D/2\leq r_m\leq D.
\]
\texttt{HamiltonianFactorFloor} retains the full one-unit
floor error and proves \(D\geq L/4\).
\texttt{HamiltonianPrimeFactorFamily} constructs one family,
with both original degree estimates retained, and proves
\[
 \left|\frac{r_m}{D}-1\right|
       \leq\frac{10B}{\sqrt L}+\frac4L.
\]
Its terminal threshold is fixed; no moving defect is
substituted into a fixed-defect eventual theorem.
\texttt{HamiltonianFactorNormalization} uses the exact identity
\[
 \frac{A_{J_m}}{r_m}-U_{2m}
 =\frac{A_{J_m}-D U_{2m}}{r_m}
       +\frac{D-r_m}{r_m}U_{2m}.
\]
\texttt{HamiltonianPrimeFactorNormalization} consumes both
terms for that same chosen family and proves
\[
 \left\|\frac{A_{J_m}}{r_m}-U_{2m}\right\|
                 \leq\frac{62B}{\sqrt L}+\frac{24}L.
\]
The shared \texttt{HamiltonianSupportMixture} theorem proves
that a positive mixture of a stochastic factor with a
full-support stochastic reference retains every original
edge, with perturbation at most twice the mixing weight.
Twenty-one core declarations in these five owners pass
focused foundational-axiom checks.

Theorem~\ref{thm:nearlyoptimalfactor} now has its complete
maintained consumer. \texttt{HamiltonianPrimeFactorReference}
proves the normalized factor's actual stochastic law and
constructs one full-support reference family
\[
 P_m=(1-\delta_m)\frac{A_{J_m}}{r_m}+\delta_m R_m,
 \qquad \delta_m=\frac1{m+1},
\]
using one fixed-terminal full-support stochastic family \(R_m\).
It retains exact symmetry, nonnegativity, every row and column
sum, and every original prime edge, with
\[
 \left\|\frac{A_{J_m}}{r_m}-P_m\right\|\leq\frac2{m+1},
 \qquad
 \|P_m-U_{2m}\|\leq
       \frac{62B}{\sqrt L}+\frac{24}L+\frac2{m+1}.
\]
\texttt{HamiltonianPrimeFactorLimits} proves the complete
terminal epsilon bounds; the same chosen family works for
every precision. \texttt{HamiltonianBipartiteEigenvalues}
derives both vanishing block sums from the actual eigenvector
equations whenever the normalized eigenvalue differs from
\(1,-1\). Thus the constant/parity reference acts as zero,
without an added orthogonality hypothesis.
\texttt{HamiltonianPrimeFactorEigenvalues} transfers this
argument through the exact vertex bijection and applies it
to the native real Hermitian eigenvalues of \(A_{J_m}\).

Finally, \texttt{HamiltonianPrimeFactorHeadline} proves,
for every \(\eta>0\), one terminal threshold for which the
same \(J_m,P_m\) satisfy
\[
 \left|\frac{r_m}{D}-1\right|\leq\eta,\qquad
 \left\|\frac{A_{J_m}}{r_m}-P_m\right\|\leq\eta,\qquad
 \left\|\frac{A_{J_m}}{r_m}-P_{0,2m}\right\|
       \leq\frac12+\eta .
\]
Every native adjacency eigenvalue unequal to \(r_m,-r_m\)
then has modulus at most \((1/2+\eta)r_m\).
The exact constant and parity eigenvector identities are
also proved. Uniformly over every other regular factor
\(K\subseteq G_{2m}\) of degree \(s\), the same terminal
interval gives \(s/D\leq1+\eta\), proving asymptotic degree
optimality. All twenty-nine core declarations in these five
new owners pass focused foundational-axiom checks.
The resilience theorem below now has its complete signed-count
consumer. The cycle-packing theorem retains its separate
formalization obligation.

\begin{theorem}[A positive proportion of the prime edges in Hamilton cycles]
\label{thm:manycycles}
For every fixed \(\kappa<1/8\) with \(\kappa>0\), and every
sufficiently large even \(N\), \(G_N\) contains at least
\[
 \left\lfloor\kappa\,\frac N{\log N}\right\rfloor
 \tag{30}
\]
pairwise edge-disjoint Hamilton cycles.
Consequently the supremal number of such cycles, divided by
\(N/\log N\), has lower limit at least \(1/8\).
\end{theorem}
\begin{proof}
Start with \(J_N\) from (29), of degree \(r=r_N\).
After deleting \(j\) edge-disjoint Hamilton cycles, the remaining
graph \(J^{(j)}\) is simple, bipartite and \((r-2j)\)-regular.
The deleted adjacency has norm at most \(2j\). Weyl's inequality,
including this full signed subtraction, gives
\[
 \lambda_2(J^{(j)})\leq(1/2+o(1))r+2j,
\]
and hence
\[
 (r-2j)-\lambda_2(J^{(j)})
                   \geq(1/2-o(1))r-4j.                    \tag{31}
\]
Fix any \(\kappa'\) strictly between \(\kappa\) and \(1/8\).
For all \(j\leq\kappa'r\), (31) leaves a fixed positive
spectral gap as a proportion of the remaining degree.
That degree is of order \(N/\log N\), so exceeds the
polylogarithmic requirement in (2), uniformly in these \(j\),
for all sufficiently large \(N\).
Equation (2) supplies another Hamilton cycle at each step.
Since \(r/(N/\log N)\to1\), this produces (30).
No independence of the successive cycles is used.
\end{proof}

\paragraph{Formalized retained-mode input for cycle packing.}
The full uniform packing conclusion above remains a formalization
obligation. Its arithmetic finite-mode input is now proved.
Let \(Q_n\) be the proved moving prime-operator cutoff, fix a
positive integer \(K\), and put \(H=27K^4\). Index the retained
modes by
\[
 {\cal U}_H=\{(q,a):1\leq q\leq H,\quad
                          a\in(\mathbb Z/q\mathbb Z)^\times\}.
\]
The exact fixed-index model \(F_{n,R,K}\) retains the coefficient
\(\mu(q)/\varphi(q)\) when \(q\leq Q_n\), and zero otherwise,
together with the original sampled character phase. The
auxiliary grid has \(R>2n\), \(2n\mid R\), and \(q\mid R\)
for every \(1\leq q\leq Q_n\).

\texttt{HamiltonianRationalTailNorm} consumes the already proved
uniform inequality \(2q^3\leq27\varphi(q)^4\), deriving
\[
 q>27K^4\quad\Longrightarrow\quad
 |\mu(q)/\varphi(q)|\leq1/K.
\]
It proves the full two-frame remainder bound
\[
 \|\mathsf R_{n,R,Q}-\mathsf R_{n,R,Q}^{\leq H}\|
 \leq K^{-1}(1+4Q^4/n).
\]
\texttt{HamiltonianRationalCutoffTransport} exchanges the moving
and fixed denominator indices exactly, retaining every signed
coefficient, phase and inactive denominator.
\texttt{HamiltonianPrimeRationalApproximation} exposes the
stronger part of the complete major/minor calculation before
its two-mode relaxation: uniformly on a terminal interval,
\[
 \|\mathsf M_n-\mathsf R_{n,R,Q_n}\|+2Q_n^4/n
                                                    \leq\varepsilon,
\]
where \(\mathsf M_n\) is the actual Mangoldt sum matrix.
Thus \texttt{HamiltonianPrimeFiniteModeApproximation} proves
\[
 \|\mathsf M_n-F_{n,R,K}\|
       \leq\varepsilon+K^{-1}(1+2\varepsilon).
\]
For every \(\eta>0\), one fixed \(K\) and one initial order
therefore give error at most \(\eta\) for every subsequent
order and every admissible grid.
The same owner proves \(F_{n,R,K}=A_nB_n\) through the
fixed vector space \(\mathbb C^{{\cal U}_H}\), with
\(|{\cal U}_H|\leq H^2\).
All seventeen core declarations in these four owners pass
focused foundational-axiom checks.
Transfer to the normalized simple prime factors, the retained
finite-mode determinant comparison and the successive
edge-disjoint cycle construction remain to be formalized.
This completed arithmetic input is not a completed packing
certificate.

\begin{theorem}[Explicit positive local resilience]
\label{thm:resilience}
For every sufficiently large even \(N\), if an arbitrary set
of edges \(F\) of \(G_N\) has
\[
 \max_v\deg_F(v)\leq \frac1{1000}\frac N{\log N},             \tag{32}
\]
then \(G_N-F\) is Hamiltonian.
\end{theorem}
\begin{proof}
Write \(D=2m/\log(2m)\), \(N=2m\), and choose the order
large enough that the ordinary degrees lie in
\([0.999D,1.001D]\), the full cut bound (28) holds, and
\[
 \|U_{2m}-P_{0,2m}\|\leq0.501,\qquad
 \|U_{2m}\|\leq1.001.
\]
Let \(Q\) be the deleted edges in the canonical odd/even
cross block. Every row and column of \(Q\) has at most
\(D/1000\) edges.

Here is the complete deletion accounting. For arbitrary
\(S,T\) in the two parts, write \(a_E,b_E\) for the two
complementary crossing counts in edge set \(E\).
Their exact finite identities are
\[
 a_G=a_{G\setminus Q}+a_Q,\qquad
 b_G=b_{G\setminus Q}+b_Q.
\]
Counting the deleted boundary from either side gives
\[
 a_Q+b_Q\leq\frac D{1000}
 \min\{|S|+|T|,\ |S^c|+|T^c|\}.
\]
Thus the remaining degrees lie in
\([0.998D,1.001D]\subseteq[0.998D,1.002D]\), and the
remaining full cut coefficient is at least \(0.199\).
Lemma~\ref{lem:regularfactor} applies with
\(\varepsilon=0.002\), \(g=0.199\).
For \(D\geq1000\), it permits
\[
 r=\lfloor(39/40)D\rfloor,\qquad 0.97D\leq r\leq D,
 \qquad 39/40\leq195/199=1-\frac{2(0.002)}{0.199}.
\]
Obtain a simple \(r\)-regular factor \(H\subseteq G_N-Q\).
Measure the full discarded adjacency from the original graph:
\[
 \|A_{G_N}-A_H\|\leq1.001D-r.
\]
This includes \(Q\) once. No independent additional loss
for \(Q\) is charged.

Both normalization terms are retained:
\[
 \frac{A_H}{r}-U_{2m}
 =\frac{A_H-DU_{2m}}r+\frac{D-r}rU_{2m}.
\]
Consequently
\[
 \left\|\frac{A_H}r-U_{2m}\right\|
 \leq \frac{1.001D-r+1.001(D-r)}r
 =2.002\,\frac Dr-2.001<0.064,
\]
and hence
\[
 \left\|\frac{A_H}r-P_{0,2m}\right\|
       \leq0.565<3/5.
\]
The cross block \(T\) of \(A_H/r\) is doubly stochastic.
Put \(a=(3/5)\log(2m)\). Its entries satisfy
\[
 0\leq T_{ij}\leq1/r\leq a/m,\qquad
 a-1\leq(3/5)\log m.
\]
The exact centered cross norm equals the preceding physical
norm. The \(3/5\) specialization of
Theorem~\ref{thm:criticalrelative} therefore gives a relative
Hamilton-cycle error at most \(m^{-3/5}<1\).
The resulting cycle weight is strictly positive.
A positive term in the finite permutation sum supplies a
one-orbit permutation, and every edge of that permutation
belongs to \(H\). Through the exact odd/even vertex
bijection this is a Hamilton cycle on the original integers,
using no deleted edge. All thresholds precede the choice of
\(Q\), so the conclusion is uniform over (32).
\end{proof}

The complete resilience consumer is formalized in
\texttt{HamiltonianPrimeResilience}.
\texttt{HamiltonianFactorDeletion} proves the exact finite
partition of every crossing count, the two-sided deleted
boundary bound, and extraction from the remaining graph,
while measuring the total adjacency loss from the original
graph. \texttt{HamiltonianResilientNormalization} proves
the floor constraints and the full \(0.064\) budget.
\texttt{HamiltonianRegularFactorKernels} constructs the
actual regular cross kernel, proves its exact stochastic
law and physical centered-norm identity, and extracts an
actual one-orbit cycle from a positive signed count.
The final theorem keeps every original distinct-integer
prime-sum edge condition and proves that no deleted edge is
used. All nineteen core declarations in these four owners
pass focused foundational-axiom checks. This proof consumes
the signed counting theorem directly; the cycle-packing
theorem still has its own formalization obligation.
